% Chapter 2 of the BARCS student guide: the calculations.
\ifstandalone
%=====================================================================
\section{What you need from Chapter 1}
%=====================================================================

This chapter continues the example of Chapter~1. Here is a short summary.

\textbf{The data.} One guide, four libraries:

\begin{center}
\begin{tabular}{lcrr}
\toprule
Library & Group & Guide count $K$ & Total count $N$\\
\midrule
1 & control & 180 & 1{,}000{,}000\\
2 & control & 240 & 1{,}200{,}000\\
3 & drug    & 120 & 1{,}000{,}000\\
4 & drug    &  80 &   800{,}000\\
\bottomrule
\end{tabular}
\end{center}

\begin{itemize}
\item \textbf{Proportion} $p = K/N$. Combined: 191 per million (control) and
  111 per million (drug). The fold change is 0.58.
\item \textbf{Log-odds} $= \ln\bigl(p/(1-p)\bigr)$: $-8.56$ (control) and
  $-9.10$ (drug).
\item \textbf{Noise.} The variance of a count is
  $N \times \mu \times (1-\mu) \times \bigl(1 + (N-1) \times \rho\bigr)$.
  Here $\mu$ is the expected proportion. $\rho$ is the correlation between
  two reads from the same library; it measures the biological noise. The
  wobble of the true proportion between replicates is
  $V = \rho \times \mu \times (1-\mu)$.
\item \textbf{Two rules.} Variances of independent values add up. Multiplying
  a value by $c$ multiplies its variance by $c \times c$.
\item \textbf{Design matrix} $X$: one row for each library, one column for
  each $\beta$. Our $X$ has the rows $(1,0)$, $(1,0)$, $(1,1)$, $(1,1)$.
\item \textbf{Prediction.} Multiply a row of $X$ by the $\beta$ values and add.
  The result is the predicted log-odds $\eta$. Then
  $\mu = e^{\eta}/(1+e^{\eta})$.
\item $\beta_0$ is the control log-odds. $\beta_1$ is the effect of the drug.
\end{itemize}
\fi

%=====================================================================
\section{Step 6: How BARCS finds the best \texorpdfstring{$\beta$}{beta} values}
%=====================================================================

BARCS cannot find $\beta_0$ and $\beta_1$ in one calculation. It uses a loop
of ``guess, check, improve''. The name of this method is
\textbf{iteratively reweighted least squares} (IRLS). Each pass through the
loop is an \textbf{iteration}.

\subsection{The loop}

\begin{enumerate}
\item \textbf{Guess.} Start with the simple log-odds of each group
  (Step~3 in Chapter~1).
\item \textbf{Predict.} Use the design matrix to predict $\eta$ and $\mu$
  for each library (Step~5 in Chapter~1).
\item \textbf{Measure the noise.} Find $\rho$ from the differences between
  the real and the expected counts (Step~7).
\item \textbf{Weigh.} Give each library a \textbf{weight}. A library with
  less noise gets a larger weight.
\item \textbf{Correct.} Calculate a corrected value, the \textbf{working
  value}, for each library.
\item \textbf{Improve.} Take a weighted average of the working values in each
  group. These averages give the new $\beta$ values.
\item \textbf{Repeat} from part~2 until the $\beta$ values stop changing.
\end{enumerate}

\subsection{The working value: a straight-line step}

BARCS works on the log-odds scale. The real log-odds of library $i$ is
$\ln\bigl((K_i/N_i) / (1 - K_i/N_i)\bigr)$. But BARCS does not use this value
directly. If $K_i = 0$, the value is $\ln(0)$, which does not exist. So BARCS
uses a straight line instead.

\paragraph{The slope of the log-odds.} Near a proportion $\mu$, the
log-odds changes by about $1 \div \bigl(\mu \times (1-\mu)\bigr)$ for each
unit of change in the proportion. Here is why. The log-odds is
$\ln(p) - \ln(1-p)$. Look at the two parts one at a time, with $p = 0.2$
and a small step of 0.001:

\begin{center}
\small
\renewcommand{\arraystretch}{1.3}
\begin{tabular}{@{}l l l l@{}}
\toprule
Part & At $p = 0.2$ & At $p = 0.201$ & Change $\div$ 0.001\\
\midrule
$\ln(p)$ & $-1.60944$ & $-1.60445$ & $4.99 \approx 1 \div 0.2 = 5$\\
$-\ln(1-p)$ & $+0.22314$ & $+0.22439$ & $1.25 = 1 \div 0.8$\\
log-odds (sum) & $-1.38629$ & $-1.38006$ & $6.24 \approx 6.25$\\
\bottomrule
\end{tabular}
\end{center}

So the slope is $1 \div p + 1 \div (1-p)$. Put the two fractions over one
bottom:
\[
\frac{1}{p} + \frac{1}{1-p} = \frac{(1-p) + p}{p \times (1-p)}
 = \frac{1}{p \times (1-p)} .
\]
Check: $1 \div (0.2 \times 0.8) = 1 \div 0.16 = 6.25$.

\paragraph{The working value.} Start at the current prediction $\eta_i$
(the log-odds of $\mu_i$). Move along the straight line by the distance
$K_i/N_i - \mu_i$:
\[
v_i = \eta_i + (K_i/N_i - \mu_i) \times \frac{1}{\mu_i \times (1-\mu_i)} .
\]
The letter $v$ stands for ``working value''. It is close to the real
log-odds, and it exists even when $K_i = 0$.

\begin{worked}[ 10a: working value of library 1]
Use the first guess $\mu_1 = 0.000190909$, so $\eta_1 = -8.563523$.
\begin{itemize}[leftmargin=1.2em]
\item Distance: $K_1/N_1 - \mu_1 = 0.000180 - 0.000190909 = -0.0000109$.
\item Slope: $1 \div (0.000190909 \times 0.999809) = 1 \div 0.000190873 =
  5{,}239$.
\item Step: $-0.0000109 \times 5{,}239 = -0.0572$.
\item Working value: $v_1 = -8.5635 + (-0.0572) = -8.6207$.
\end{itemize}
The real log-odds is $\ln(0.000180 / 0.999820) = -8.6224$. The straight line
gives $-8.6207$. They are very close.
\end{worked}

\subsection{The weight: one divided by the variance}

\paragraph{Why weights?} Suppose you have two measurements of the same
thing. One has the variance 1. The other has the variance 4 (it is noisier).
\begin{itemize}
\item With equal weights, the average has the variance
  $(1 + 4) \div (2 \times 2) = 1.25$.
\item With the weights $1 \div 1 = 1$ and $1 \div 4 = 0.25$, the average has
  the variance $1 \div (1 + 0.25) = 0.8$. That is smaller.
\end{itemize}
The rule ``weight $= 1 \div$ variance'' gives the average with the smallest
variance. So BARCS gives each working value the weight
\[
w_i = \frac{1}{\text{variance of } v_i} .
\]

\paragraph{The variance of the working value.} We build it from Chapter~1
in three steps. The right column uses library 1 in iteration 1
($N = 1{,}000{,}000$, $\mu = 0.000190909$, $\rho = 0.0000003841$).

\begin{center}
\small
\renewcommand{\arraystretch}{1.45}\setlength{\tabcolsep}{4pt}
\begin{tabular}{@{}r@{\;}l l l@{}}
\toprule
& \textbf{Quantity} & \textbf{Formula} & \textbf{Library 1}\\
\midrule
1. & variance of $K_i$ & $N_i \times \mu_i \times (1-\mu_i) \times \bigl(1 + (N_i-1) \times \rho\bigr)$ & $190.87 \times 1.3841 = 264.19$\\
2. & variance of $K_i / N_i$ & line 1 $\div (N_i \times N_i)$ & $264.19 \div 10^{12}$\\ & & & $= 2.642 \times 10^{-10}$\\
3. & variance of $v_i$ & line 2 $\div \bigl(\mu_i \times (1-\mu_i)\bigr)^2$ & $2.642 \times 10^{-10} \div 0.000190873^2$\\ & & & $= 0.007251$\\
4. & weight $w_i$ & $1 \div$ line 3 & $1 \div 0.007251 = 137.90$\\
\bottomrule
\end{tabular}
\end{center}

\begin{itemize}
\item \textbf{Line 1} is the variance formula of Chapter~1, Step~4.
\item \textbf{Line 2:} dividing a value by $N$ divides its variance by
  $N \times N$ (Chapter~1, Part~4 used the same rule with multiplying).
\item \textbf{Line 3:} the working value multiplies $K_i/N_i - \mu_i$ by the
  slope $1 \div \bigl(\mu_i \times (1-\mu_i)\bigr)$. So the variance is
  multiplied by the slope squared. Adding $\eta_i$ does not change the
  variance.
\item \textbf{Line 4:} put lines 1--3 together and simplify. One
  $N_i$ and one $\mu_i \times (1-\mu_i)$ cancel:
\end{itemize}
\[
w_i = \frac{1}{\text{variance of } v_i}
 = \frac{N_i \times N_i \times \bigl(\mu_i \times (1-\mu_i)\bigr)^2}
        {N_i \times \mu_i \times (1-\mu_i) \times \bigl(1 + (N_i-1) \times \rho\bigr)}
 = \frac{N_i \times \mu_i \times (1-\mu_i)}{1 + (N_i - 1) \times \rho} .
\]
Check: $190.87 \div 1.3841 = 137.90$, the same as line 4.

A library with more reads gets a bigger top, so a bigger weight. More
biological noise ($\rho$) gives a bigger bottom, so a smaller weight.

\subsection{The new \texorpdfstring{$\beta$}{beta} values}

For our design matrix, the new $\beta$ values are weighted averages of the
working values:
\[
\beta_0 = \frac{w_1 \times v_1 + w_2 \times v_2}{w_1 + w_2}, \qquad
\beta_0 + \beta_1 = \frac{w_3 \times v_3 + w_4 \times v_4}{w_3 + w_4}.
\]
(For other design matrices, BARCS solves a small set of equations instead.
The idea is the same.)

\paragraph{Why a loop?} The weights and the working values use $\mu$. But
$\mu$ comes from the $\beta$ values, and the $\beta$ values come from the
weights. So BARCS starts with a guess and repeats until nothing changes.

\subsection{Iteration 1 in full}

\begin{worked}[ 10b: iteration 1]
\textbf{Guess.} From Step~3 in Chapter~1, the control log-odds is $-8.563522$ and the
drug log-odds is $-9.104869$. So $\beta_0 = -8.563522$ and
$\beta_1 = -9.104869 - (-8.563522) = -0.541347$.

\textbf{Predict.} The control libraries get $\mu = 0.000190909$. The drug
libraries get $\mu = 0.000111111$.

\textbf{Measure the noise.} Step~7 gives $\rho = 0.0000003841$ for these
predictions.

\textbf{Weigh library 1.}
Top: $1{,}000{,}000 \times 0.000190909 \times 0.999809 = 190.87$.\\
Bottom: $1 + 999{,}999 \times 0.0000003841 = 1.3841$.\\
Weight: $w_1 = 190.87 \div 1.3841 = 137.90$.

\textbf{Correct library 1.} Example 10a: $v_1 = -8.6207$.

The same calculations for all four libraries give:

\begin{tabular}{@{}lrrr@{}}
\toprule
Library & Weight $w_i$ & Working value $v_i$ & $w_i \times v_i$\\
\midrule
1 (control) & 137.905 & $-8.62068$ & $-1188.84$\\
2 (control) & 156.785 & $-8.51589$ & $-1335.16$\\
3 (drug)    &  80.269 & $-9.02486$ &  $-724.42$\\
4 (drug)    &  67.988 & $-9.20488$ &  $-625.82$\\
\bottomrule
\end{tabular}

\textbf{Improve.}\\
Control: $\beta_0 = (-1188.84 - 1335.16) \div (137.905 + 156.785)
 = -2524.00 \div 294.69 = -8.564929$.\\
Drug: $\beta_0 + \beta_1 = (-724.42 - 625.82) \div (80.269 + 67.988)
 = -1350.24 \div 148.257 = -9.107414$.\\
So $\beta_1 = -9.107414 - (-8.564929) = -0.542486$.
\end{worked}

\subsection{The next iterations}

BARCS repeats the loop with the new values. The changes become very small:

\begin{center}
\begin{tabular}{crr}
\toprule
After iteration & $\beta_0$ & $\beta_1$\\
\midrule
0 (the guess) & $-8.563522$ & $-0.541347$\\
1 & $-8.564929$ & $-0.542486$\\
2 & $-8.564934$ & $-0.542492$\\
3 & $-8.564934$ & $-0.542492$\\
\bottomrule
\end{tabular}
\end{center}

After iteration 3, the values do not change in the first six decimal places.
BARCS stops. The final values are $\beta_0 = -8.565$ and
$\beta_1 = -0.542$.

\begin{remember}
The weighted average gives a slightly different answer from the simple
log-odds of Step~3 in Chapter~1 ($-0.5425$ instead of $-0.5413$). The difference comes
from the weights. Library~2 has the most reads and the least relative noise.
So it counts the most.
\end{remember}

%=====================================================================
\section{Step 7: How BARCS measures the biological noise}
%=====================================================================

Step~4 in Chapter~1 used $\rho = 0.000000386$. This section shows where the number comes
from.

\subsection{Pearson residuals}

For each library, BARCS subtracts the expected count from the real count.
It divides the result by the standard deviation from Step~4 in Chapter~1. The result is a
\textbf{Pearson residual}:
\[
\text{Pearson residual}_i =
\frac{K_i - N_i \times \mu_i}{\sqrt{N_i \times \mu_i \times (1-\mu_i) \times \bigl(1 + (N_i-1) \times \rho\bigr)}}.
\]
A Pearson residual tells how many standard deviations a library is from its
prediction.

\subsection{The rule for \texorpdfstring{$\rho$}{rho}}

\paragraph{Why the sum should be the degrees of freedom.} A Pearson
residual is a distance measured in standard deviations. If the noise model
is correct, its square is 1 on average: (squared distance) $\div$ (variance)
$= $ variance $\div$ variance $= 1$. So with four libraries, the four
squares add up to about 4.

But BARCS does not know the true $\mu$. It uses the fitted $\mu$, which the
$\beta$ values pull toward the data. The fit makes the residuals a little
smaller. Each fitted $\beta$ value takes away 1 from the sum, on average. So
the expected sum is
\[
(\text{number of libraries}) - (\text{number of } \beta \text{ values})
 = 4 - 2 = 2 = \text{df}.
\]
A computer checked this. It made 20{,}000 artificial guides with the true
$\rho$ and four libraries like ours. With the true $\mu$, the sum of the
squared Pearson residuals was 4.00 on average. With the fitted $\mu$, it was
2.00 on average.

\paragraph{The rule.} BARCS turns this around. It chooses $\rho$ so that the
sum is exactly the degrees of freedom:
\[
\text{sum of the squared Pearson residuals} = \text{df}.
\]
(Statisticians call this a \emph{method of moments}: make a sum that you can
calculate equal to the value it should have on average.)

\begin{worked}[ 11: start with \texorpdfstring{$\rho = 0$}{rho = 0}]
Use the expected counts from Step~5 in Chapter~1. With $\rho = 0$, the bottom of each
residual is the counting noise only.

\begin{tabular}{@{}lrrr@{}}
\toprule
Library & $(K - N \times \mu)^2$ & $N \times \mu \times (1-\mu)$ & Squared residual\\
\midrule
1 & $(-10.64)^2 = 113.21$ & 190.60 & 0.594\\
2 & $(+11.23)^2 = 126.16$ & 228.72 & 0.552\\
3 & $(+9.17)^2 = 84.14$   & 110.82 & 0.759\\
4 & $(-8.66)^2 = 75.03$   &  88.65 & 0.846\\
\midrule
  &  & Sum & 2.751\\
\bottomrule
\end{tabular}

\smallskip
The sum is 2.751. The degrees of freedom are 2. The sum is too big. So
there is more noise than counting noise alone. We need $\rho > 0$.
\end{worked}

\begin{worked}[ 12: guess and check for \texorpdfstring{$\rho$}{rho}]
A larger $\rho$ makes the bottom of each residual larger. So the sum goes
down. BARCS tries values until the sum equals 2:

\begin{tabular}{@{}rr@{}}
\toprule
$\rho$ & Sum of squared residuals\\
\midrule
0           & 2.751\\
0.0000001   & 2.506\\
0.0000002   & 2.302\\
0.0000003   & 2.129\\
0.000000386 & 2.000\\
0.0000005   & 1.851\\
\bottomrule
\end{tabular}

\smallskip
The answer is $\rho = 0.000000386$. If the sum is already 2 or less at
$\rho = 0$, BARCS uses $\rho = 0$.
\end{worked}

The ratio $2.751 \div 2 = 1.376$ is also useful. It is the
\textbf{variance inflation}. At $\rho = 0$, the sum should be about 2; it is
2.751. So the variance is about 1.38 times bigger than counting noise alone.
This is close to the extra factor $1 + (N-1) \times \rho = 1.386$ of
Chapter~1. Step~9 uses the inflation.

%=====================================================================
\section{Step 8: Is the effect real?}
%=====================================================================

The effect is $\beta_1 = -0.542$. But replicates are noisy. Maybe the true
effect is 0 and the noise made the drop. BARCS tests this question in four
parts.

\subsection{The standard error}

The \textbf{standard error} (SE) tells how much the estimate of $\beta_1$
can change because of noise. A small SE means a precise estimate.

For our design matrix, the standard error comes from the weights of Step~6.
Add the weights of each group:
\[
W_{\text{control}} = w_1 + w_2, \qquad W_{\text{drug}} = w_3 + w_4.
\]
Then
\[
\text{SE}(\beta_1) = \sqrt{\frac{1}{W_{\text{control}}} +
\frac{1}{W_{\text{drug}}}}.
\]
\paragraph{Where this formula comes from.} It uses only three facts:
the weight is $1 \div$ variance (Step~6); variances of independent values
add up; and multiplying a value by $c$ multiplies its variance by
$c \times c$.

\begin{center}
\small
\renewcommand{\arraystretch}{1.45}\setlength{\tabcolsep}{4pt}
\begin{tabular}{@{}r@{\;}l l@{}}
\toprule
& \textbf{Step} & \textbf{Why}\\
\midrule
1. & variance of $v_i = 1 \div w_i$ & definition of the weight\\
2. & $\beta_0 = \dfrac{w_1}{W_c} \times v_1 + \dfrac{w_2}{W_c} \times v_2$ \quad ($W_c = w_1 + w_2$) & Step~6\\
3. & variance of $\beta_0 = \Bigl(\dfrac{w_1}{W_c}\Bigr)^2 \times \dfrac{1}{w_1} + \Bigl(\dfrac{w_2}{W_c}\Bigr)^2 \times \dfrac{1}{w_2}$ & multiply by $c \Rightarrow$ variance $\times\, c \times c$; add\\
4. & $\phantom{\text{variance of }\beta_0} = \dfrac{w_1 + w_2}{W_c \times W_c} = \dfrac{1}{W_c}$ & one $w$ cancels in each part\\
5. & variance of $(\beta_0 + \beta_1) = 1 \div W_{\text{drug}}$ & the same for the drug libraries\\
6. & variance of $\beta_1 = \dfrac{1}{W_c} + \dfrac{1}{W_{\text{drug}}}$ & $\beta_1 = (\beta_0 + \beta_1) - \beta_0$; variances add\\
7. & $\text{SE}(\beta_1) = \sqrt{\text{line 6}}$ & SE is the standard deviation of the estimate\\
\bottomrule
\end{tabular}
\end{center}

Line~6 adds the variances even though it subtracts the values. The reason:
$-1 \times -1 = +1$, so a minus sign does not change a variance.

More weight means less noise, so the SE becomes smaller. (For other design
matrices, BARCS uses the inverse of the matrix $X^{\top} W X$. For our design,
that calculation gives exactly this formula.)

\begin{worked}[ 13: the standard error]
The final weights are 137.524, 156.322, 79.955 and 67.737.

$W_{\text{control}} = 137.524 + 156.322 = 293.846$.\\
$W_{\text{drug}} = 79.955 + 67.737 = 147.692$.

$1 \div 293.846 = 0.003403$ (line 4: $(137.524 + 156.322) \div 293.846^2$ gives the same).\\
$1 \div 147.692 = 0.006771$.\\
Sum: $0.003403 + 0.006771 = 0.010174$.\\
$\text{SE} = \sqrt{0.010174} = 0.101$.

The control group has about twice the weight of the drug group. So most of
the uncertainty comes from the drug group.
\end{worked}

\subsection{The t value}

The \textbf{t value} compares the effect with its standard error:
\[
t = \frac{\beta_1}{\text{SE}}.
\]
A $t$ value far from 0 means that the effect is large compared with the
noise.

\subsection{Degrees of freedom}

The \textbf{degrees of freedom} (df) tell how much information we have
about the noise:
\[
\text{df} = (\text{rows of } X) - (\text{columns of } X)
 = (\text{number of libraries}) - (\text{number of } \beta \text{ values}).
\]
Fewer degrees of freedom mean less information. Then BARCS needs a bigger
$t$ value before it accepts an effect as real. (The reason is that the SE
itself is estimated from few libraries, so it can be too small by chance.
W.\,S. Gosset, who wrote as ``Student'', worked out the table for this case
in 1908.\footnote{Student (1908). The probable error of a mean.
\emph{Biometrika} 6: 1--25.})

\begin{center}
\begin{tabular}{lcccc}
\toprule
Degrees of freedom & 1 & 2 & 14 & very many\\
$|t|$ needed for $p < 0.05$ & 12.71 & 4.30 & 2.14 & 1.96\\
\bottomrule
\end{tabular}
\end{center}

\subsection{The p-value}

The \textbf{p-value} answers this question: ``If the true effect is 0, how
often does noise alone give a $t$ value this far from 0?'' A small p-value
means that noise alone seldom gives such a result. Scientists often use 0.05
as a limit.

\begin{worked}[ 14: testing the effect]
\begin{enumerate}[leftmargin=1.4em]
\item $t = -0.542 \div 0.101 = -5.38$.
\item $X$ has 4 rows and 2 columns. So $\text{df} = 4 - 2 = 2$.
\item The table shows that 2 degrees of freedom need $|t| > 4.30$. Our
  $|t| = 5.38$. It is larger than 4.30. So $p < 0.05$.
\item A computer gives the exact p-value: $p = 0.033$.
\end{enumerate}
The p-value is less than 0.05. So the drop is probably not only noise.
\end{worked}

\begin{remember}
The degrees of freedom come from the number of libraries, not from the
number of reads. A library with 100 million reads still counts as one row
of $X$.
\end{remember}

%=====================================================================
\section{Step 9: Borrow information from other guides}
%=====================================================================

With only 2 degrees of freedom, the noise estimate of one guide is not
reliable. But a screen has thousands of guides. Most guides have a similar
amount of noise. So BARCS combines two estimates:

\begin{itemize}
\item the variance inflation of this guide (Step~7), and
\item the typical variance inflation of guides with similar counts.
\end{itemize}

This method is \textbf{moderation}.\footnote{BARCS uses the
empirical-Bayes method of G.\,K. Smyth (2004). Linear models and empirical
Bayes methods for assessing differential expression in microarray
experiments. \emph{Statistical Applications in Genetics and Molecular
Biology} 3: Article 3.} BARCS takes a weighted average. Each
estimate gets a number of ``votes''. The guide's own estimate gets its
degrees of freedom as votes. The typical estimate gets the \textbf{prior
degrees of freedom}, $d_0$. BARCS calculates $d_0$ and the typical value
from all the guides in the screen.

\[
\text{moderated inflation} =
\frac{\text{df} \times \text{own inflation} + d_0 \times
\text{typical inflation}}{\text{df} + d_0}.
\]

\begin{worked}[ 15: moderation with simple numbers]
A guide has 3 degrees of freedom. Its own inflation is 4.0. The typical
inflation is 1.5, with $d_0 = 12$ votes.
\[
\text{moderated inflation} = \frac{3 \times 4.0 + 12 \times 1.5}{3 + 12}
 = \frac{12 + 18}{15} = \frac{30}{15} = 2.0.
\]
The guide's value moves from 4.0 toward 1.5. It stops at 2.0.
\end{worked}

\subsection{What moderation does to the test}

BARCS uses the moderated inflation in three places:
\begin{enumerate}
\item It multiplies the standard error by
  $\sqrt{\text{moderated inflation} \div \text{own inflation}}$. Here is
  why. The variance of $\beta_1$ grows in proportion to the inflation. To
  swap the own inflation for the moderated one, multiply the variance by
  (moderated $\div$ own). The SE is the square root of the variance, so it
  is multiplied by the square root. (If the own inflation is below 1, BARCS
  uses 1, so moderation never makes a test less careful than counting noise
  alone.)
\item It recalculates $t$ with the new standard error.
\item It adds the $d_0$ votes to the degrees of freedom. So the test uses
  $\text{df} + d_0$.
\end{enumerate}

\begin{worked}[ 16: moderation for our guide]
Our guide has own inflation 1.376 (Step~7) and 2 degrees of freedom.
Suppose that the screen gives typical inflation 1.5 with $d_0 = 12$.

Moderated inflation: $(2 \times 1.376 + 12 \times 1.5) \div (2 + 12)
 = (2.751 + 18) \div 14 = 1.482$.

Factor for the SE: $\sqrt{1.482 \div 1.376} = \sqrt{1.077} = 1.038$.\\
New SE: $0.101 \times 1.038 = 0.105$.\\
New $t$: $-0.542 \div 0.105 = -5.18$.\\
New degrees of freedom: $2 + 12 = 14$.

With 14 degrees of freedom, $|t|$ must be larger than 2.14 for
$p < 0.05$. Our $|t| = 5.18$. The new p-value is 0.00014.

The SE became slightly larger, because the typical guide is a little noisier
than our guide. But the degrees of freedom grew from 2 to 14. So the test
became much more powerful.
\end{worked}

BARCS uses moderation when it has 50 or more usable guides.
%=====================================================================
\section{Step 10: Check the results with control guides}
%=====================================================================

A good screen contains \textbf{non-targeting control guides}. These guides
switch off no gene. So their true effect is 0.

If the noise model is correct, about 5\% of control guides get
$p < 0.05$ by chance. If many more control guides get $p < 0.05$, the noise
model is too small. Then BARCS makes all the tests wider.

\begin{worked}[ 17: calibration with control guides]
A screen has 100 control guides. We expect about 5 of them to get
$p < 0.05$. But 15 of them get $p < 0.05$. The noise is too small.

BARCS finds a \textbf{scale factor} that makes the control guides behave
correctly. It sorts the 100 control guides by the size of $|t|$ (the value
without its sign). It takes the value that 95 of the 100 guides stay below.
Suppose that this value is 2.47. With many degrees of freedom, the $t$ table
says that this value should be 1.96. So the scale factor is
$2.47 \div 1.96 = 1.26$. BARCS divides every $t$ value by 1.26.

A gene with $t = 3.5$ now gets $t = 3.5 \div 1.26 = 2.78$. Its p-value
becomes larger. This change protects you from false results.
\end{worked}

The scale factor is never less than 1. So calibration can only make the
results more careful.

%=====================================================================
\section{Step 11: From guides to genes}
%=====================================================================

Each gene usually has several guides, often 3 to 6. Each guide gives its
own result. BARCS combines them into one result for the gene. It uses the
\textbf{Stouffer method} in three parts.

\begin{enumerate}
\item Change each guide's p-value into a \textbf{z score}. A z score tells
  how many standard deviations a result is from 0. The sign of the z score
  follows the sign of the effect.
\item Add the z scores.
\item Divide the sum by the square root of the number of guides, $m$.
\end{enumerate}
\[
Z_{\text{gene}} = \frac{z_1 + z_2 + \cdots + z_m}{\sqrt{m}}.
\]
\paragraph{Why divide by $\sqrt{m}$?} If the gene has no effect, each $z$
score has the variance 1. The guides are independent, so the variance of the
sum is $1 + 1 + \cdots + 1 = m$ (Chapter~1, Part~2: variances add).
Dividing a value by $\sqrt{m}$ divides its variance by
$\sqrt{m} \times \sqrt{m} = m$. So $Z_{\text{gene}}$ has the variance
$m \div m = 1$ again, and the same z table works for it.\footnote{S.\,A.
Stouffer and others (1949). \emph{The American Soldier}, Vol.~1. Princeton
University Press.}

To change a p-value into a z score, use the table of the normal
distribution. A two-sided p-value of 0.05 belongs to $z = 1.96$. Some
other pairs:

\begin{center}
\begin{tabular}{lccccc}
\toprule
Two-sided p-value & 0.317 & 0.134 & 0.046 & 0.012 & 0.0027\\
$z$ score & 1.0 & 1.5 & 2.0 & 2.5 & 3.0\\
\bottomrule
\end{tabular}
\end{center}

If the effect is negative (the guide drops), the z score gets a minus sign.

\begin{worked}[ 18: combining three guides]
A gene has $m = 3$ guides. Their p-values are 0.046, 0.134 and 0.012. Their
z scores are 2.0, 1.5 and 2.5.
\[
Z_{\text{gene}} = \frac{2.0 + 1.5 + 2.5}{\sqrt{3}} = \frac{6.0}{1.732}
 = 3.46.
\]
This $Z$ gives $p = 0.00053$ for the gene. No single guide is that strong.
Together, the three guides give strong evidence.
\end{worked}

\subsection{When guides share noise}

Sometimes the guides of one gene share noise. Then they agree more than
they should by chance. BARCS measures this sharing with the
\textbf{guide correlation}, $r$. Then the variance of the sum is not $m$.
Use the squares-and-pairs rule from Chapter~1, Part~4:
\begin{itemize}
\item $m$ squares, each with the variance 1: together $m$;
\item $m \times (m-1)$ pairs, each with the covariance $r$: together
  $m \times (m-1) \times r$.
\end{itemize}
So BARCS divides by the square root of this larger variance:
\[
Z_{\text{gene}} = \frac{z_1 + z_2 + \cdots + z_m}{\sqrt{m + m \times (m-1) \times r}}.
\]

\begin{worked}[ 19: shared noise]
Use the same three guides and $r = 0.1$.
The bottom becomes $\sqrt{3 + 3 \times 2 \times 0.1} = \sqrt{3.6} = 1.897$.
So $Z_{\text{gene}} = 6.0 \div 1.897 = 3.16$, and $p = 0.0016$.
The evidence is still strong, but less strong than before.
\end{worked}

%=====================================================================
\section{Step 12: Many genes and false discoveries}
%=====================================================================

A screen tests about 20{,}000 genes. Suppose that no gene has an effect. If
we use the limit $p < 0.05$, about 5\% of the genes still pass by chance.
That is about 1{,}000 false results.

So BARCS uses the \textbf{false discovery rate} (FDR). The FDR is the
expected fraction of false results in your list of hits.

\paragraph{The idea.} When a gene has no effect, its p-value is equally
likely to be anywhere from 0 to 1. So a fraction $p$ of such genes get a
p-value of $p$ or less by chance. Example: of 20{,}000 genes with no effect,
about $20{,}000 \times 0.001 = 20$ get $p \le 0.001$. Now suppose that your
list ``all genes with $p \le 0.001$'' has 100 genes. At most about 20 of them
are there by chance. So the false fraction is about $20 \div 100 = 0.2$. In
general:
\[
\text{FDR of a list} \approx \frac{p \times (\text{number of genes})}{\text{number of genes in the list}} .
\]
If you sort the genes by p-value, the list that ends at a gene has
``rank'' genes in it. That gives part~3 below.\footnote{Y. Benjamini and Y.
Hochberg (1995). Controlling the false discovery rate: a practical and
powerful approach to multiple testing. \emph{Journal of the Royal
Statistical Society B} 57: 289--300.}

BARCS calculates it with the \textbf{Benjamini--Hochberg method}:

\begin{enumerate}
\item Sort the genes from the smallest p-value to the largest.
\item Give each gene a rank: 1, 2, 3, and so on.
\item For each gene, calculate $p \times (\text{number of genes}) \div
  \text{rank}$.
\item Go from the bottom of the list to the top. If a value is larger than
  the value below it, replace it with the value below it. (A gene is in
  every longer list too. If a longer list has a smaller FDR, the gene can
  use that list.)
\end{enumerate}

\begin{worked}[ 20: false discovery rate for five genes]
\begin{tabular}{@{}cccll@{}}
\toprule
Gene & Rank & $p$ & $p \times 5 \div \text{rank}$ & FDR (after part 4)\\
\midrule
A & 1 & 0.001 & $0.001 \times 5 \div 1 = 0.005$ & 0.005\\
B & 2 & 0.008 & $0.008 \times 5 \div 2 = 0.020$ & 0.020\\
C & 3 & 0.039 & $0.039 \times 5 \div 3 = 0.065$ & 0.051 (from D)\\
D & 4 & 0.041 & $0.041 \times 5 \div 4 = 0.051$ & 0.051\\
E & 5 & 0.600 & $0.600 \times 5 \div 5 = 0.600$ & 0.600\\
\bottomrule
\end{tabular}

\smallskip
With the limit $\text{FDR} < 0.05$, genes A and B are hits. With the limit
$\text{FDR} < 0.10$, genes A, B, C and D are hits.
\end{worked}

%=====================================================================
\section{Two special cases}
%=====================================================================

\subsection{A guide with zero reads}

Sometimes a guide has 0 reads in all control libraries. Then the log-odds
of the control group is $\ln(0)$. This value does not exist, because no
power of $e$ equals 0. Older versions of BARCS could not fit such a guide.
So they removed it.

BARCS version 0.2.1 adds a very small number to the counts first. This
number is less than one read. The method is the \textbf{Firth correction}.
Now the guide gets a value, and BARCS keeps it.

\subsection{A very large change}

Sometimes a guide grows by hundreds of times. The usual test from Step~8
can then fail. The standard error grows faster than the effect, so the
$t$ value becomes small. BARCS version 0.2.1 has a second test, the
\textbf{likelihood-ratio test}. It compares how well the model fits with
the effect and without the effect. The option \texttt{test = "auto"} uses
this test only for guides with very large changes.

%=====================================================================
\section{Summary: the full path}
%=====================================================================

\begin{center}
\begin{tikzpicture}[node distance=5mm and 4mm,
  box/.style={draw=barcsblue,fill=barcsblue!6,rounded corners=2pt,
              align=center,minimum height=9mm,text width=3.6cm,font=\small},
  arr/.style={-{Stealth},thick,barcsblue}]
\node[box] (a) {Read counts $K$\\and totals $N$};
\node[box,right=of a] (b) {Proportions\\$p = K/N$};
\node[box,right=of b] (c) {Design matrix $X$,\\fit $\beta$ by IRLS};
\node[box,below=of c] (d) {Noise: $\rho$ from\\Pearson residuals;\\moderation};
\node[box,left=of d] (e) {Test: $t$, df,\\p-value};
\node[box,left=of e] (f) {Check with\\control guides};
\node[box,below=of f] (g) {Combine guides\\into genes ($Z$)};
\node[box,right=of g] (h) {False discovery\\rate (FDR)};
\node[box,right=of h] (i) {List of\\important genes};
\draw[arr] (a)--(b); \draw[arr] (b)--(c); \draw[arr] (c)--(d);
\draw[arr] (d)--(e); \draw[arr] (e)--(f); \draw[arr] (f)--(g);
\draw[arr] (g)--(h); \draw[arr] (h)--(i);
\end{tikzpicture}
\end{center}

\subsection*{All symbols in one table}

\begin{center}
\small
\begin{tabular}{@{}l p{6.3cm} l@{}}
\toprule
Symbol & Meaning & Value in our example\\
\midrule
$K$ & read count of the guide in one library & 180, 240, 120, 80\\
$N$ & total read count of one library & 1.0, 1.2, 1.0, 0.8 million\\
$p$ & proportion, $K/N$ & 191 and 111 per million\\
fold change & $p_{\text{drug}} \div p_{\text{control}}$ & 0.58\\
$\log_2$ & number of doublings & $\log_2(0.58) = -0.78$\\
$\ln$ & natural logarithm (powers of $e = 2.718$) & $\ln(0.58) = -0.54$\\
log-odds & $\ln\bigl(p/(1-p)\bigr)$ & $-8.56$ (control)\\
$\mu$ & expected proportion & 0.000191 (control)\\
$\rho$ & correlation of two reads from one library (biological noise) & 0.000000386\\
$X$ & design matrix: a row per library, a column per $\beta$ & 4 rows, 2 columns\\
$\beta_0$ & control log-odds in the model & $-8.565$\\
$\beta_1$ & effect of the drug & $-0.542$\\
$\eta$ & predicted log-odds of a library (row of $X$ times $\beta$) & $-8.565$, $-9.107$\\
$w$ & weight of a library in the fit & 137.5 (library 1)\\
working value & prediction plus a straight-line step (Example 10a) & $-8.621$ (library 1)\\
Pearson residual & distance from the prediction in SDs & $-0.65$ (library 1)\\
inflation & noise compared with counting noise & 1.376\\
SE & standard error of the effect & 0.101\\
$t$ & effect divided by SE & $-5.38$\\
df & libraries minus $\beta$ values & $4 - 2 = 2$\\
$p$-value & chance of this result if the effect is 0 & 0.033\\
$d_0$ & prior degrees of freedom (votes) & 12 (Example 15)\\
$z$, $Z$ & z score of a guide, of a gene & 3.46 (Example 18)\\
$m$ & number of guides for a gene & 3\\
$r$ & guide correlation (shared noise) & 0.1 (Example 19)\\
FDR & expected fraction of false hits & 0.005 for gene A\\
\bottomrule
\end{tabular}
\end{center}

Note: the letter $p$ has two uses in this guide. Alone, $p$ is a proportion.
In the word ``p-value'', it is a probability. The text always tells you
which one it is.


\newpage
\input{\guideroot chapter2/exercises_ch2}
